% !TEX root = master.tex
% !TEX program = lualatex

\chapter{Pointeurs - 3}
\label{chap:pointeurs_}

\section{Chaînes de caractères en C : principes fondamentaux}

\subsection{1. Différence fondamentale avec Java}

En C :

\[
\boxed{\text{Il n’existe pas de type string.}}
\]

Une « chaîne de caractères » est simplement :

\[
\boxed{\text{un tableau de } char \text{ terminé par } '\textbackslash 0'}
\]

Il n’y a aucun objet spécial.

\subsection{2. Convention de terminaison}

La longueur n’est pas stockée explicitement.
Elle est indiquée par une convention :

\[
\boxed{\text{le caractère nul } '\textbackslash 0'}
\]

Exemple en mémoire :

\[
H \quad E \quad L \quad L \quad O \quad '\textbackslash 0'
\]

C’est uniquement la présence du caractère nul qui marque la fin.

\subsection{3. Littéraux de chaînes}

\begin{lstlisting}[language=C]
"Bonjour"
\end{lstlisting}

est équivalent en mémoire à :

\begin{lstlisting}[language=C]
{'B','o','n','j','o','u','r','\0'}
\end{lstlisting}

Remarque :

\begin{itemize}
\item Les guillemets simples `'A'` → un seul caractère.
\item Les guillemets doubles `"ABC"` → tableau de caractères + `'\0'`.
\end{itemize}

Si on veut écrire un guillemet dans la chaîne :

\begin{lstlisting}[language=C]
"Il a dit \"Bonjour\""
\end{lstlisting}

\subsection{4. Initialisation d’un tableau}

\begin{lstlisting}[language=C]
char s[] = "Bonjour";
\end{lstlisting}

Ici :

\begin{itemize}
\item s est un tableau.
\item Il y a copie automatique des caractères.
\item La taille inclut le caractère nul.
\end{itemize}

C’est équivalent à :

\begin{lstlisting}[language=C]
char s[] = {'B','o','n','j','o','u','r','\0'};
\end{lstlisting}

\subsection{5. Pointeur n’est PAS string}

\begin{lstlisting}[language=C]
char* s;
\end{lstlisting}

Ce n’est PAS une string.

C’est simplement :

\[
\boxed{\text{un pointeur vers char}}
\]

À ce stade, il pointe vers n’importe quoi.
Il doit être :

\begin{itemize}
\item soit initialisé à \texttt{NULL},
\item soit alloué dynamiquement.
\end{itemize}

\subsection{6. Allocation dynamique correcte}

Toujours allouer :

\[
n + 1
\]

pour inclure le caractère nul.

\begin{lstlisting}[language=C]
size_t n = 10;
char* s = calloc(n + 1, sizeof(char));
\end{lstlisting}

Pourquoi \texttt{n + 1} ?

Parce que nous voulons manipuler n caractères utiles,
et garder une place pour '\textbackslash 0'.

\subsection{7. Erreur classique}

\begin{lstlisting}[language=C]
char* s = "Bonjour";   // Mauvaise pratique
\end{lstlisting}

Pourquoi ?

Parce que :

\begin{itemize}
\item "Bonjour" est un littéral constant.
\item Il est stocké en mémoire en lecture seule.
\item s pointe vers une zone que vous n’avez pas le droit de modifier.
\end{itemize}

Si vous faites :

\begin{lstlisting}[language=C]
s[0] = 'A';   // comportement indefini
\end{lstlisting}

Vous risquez un crash.

\subsection{8. Cas correct 1 : chaîne constante}

Si vous ne modifiez jamais la chaîne :

\begin{lstlisting}[language=C]
const char* s = "Bonjour";
\end{lstlisting}

Ici :

\begin{itemize}
\item s pointe vers une zone constante.
\item Toute tentative de modification est interdite par le compilateur.
\end{itemize}

\subsection{9. Cas correct 2 : chaîne modifiable}

Si vous voulez modifier la chaîne :

\begin{lstlisting}[language=C]
char* s = calloc(256, sizeof(char));
strcpy(s, "Bonjour");
\end{lstlisting}

Ici :

\begin{itemize}
\item La zone est allouée dynamiquement.
\item On copie les caractères (y compris '\textbackslash 0').
\item On peut modifier le contenu.
\end{itemize}

Exemple :

\begin{lstlisting}[language=C]
s[0] = 'A';
\end{lstlisting}

La chaîne devient :

\[
"Aonjour"
\]

\subsection{10. Différence clé : copie vs pointeur}

\paragraph{Initialisation de tableau}

\begin{lstlisting}[language=C]
char s[] = "Bonjour";
\end{lstlisting}

→ Copie réelle des caractères.

\paragraph{Initialisation de pointeur}

\begin{lstlisting}[language=C]
char* s = "Bonjour";
\end{lstlisting}

→ Copie d’adresse uniquement.

Il n’y a aucune copie du contenu.

\subsection{11. Règle essentielle}

\[
\boxed{
\text{En C, une chaîne est juste un tableau de char terminé par '\textbackslash 0'}
}
\]

Il n’existe aucun objet string.

Tout repose sur :

\begin{itemize}
\item la gestion explicite de la mémoire,
\item le caractère nul,
\item la distinction tableau vs pointeur.
\end{itemize}

\section{Fonctions utilitaires pour les chaînes de caractères}

Heureusement, la bibliothèque standard fournit de nombreuses fonctions
pour manipuler les chaînes de caractères.

Toutes se trouvent dans :

\begin{lstlisting}[language=C]
#include <string.h>
\end{lstlisting}

Nous ne détaillons pas toutes les fonctions ici.
Il est essentiel de consulter les man pages :

\[
\boxed{\texttt{man 3 string}}
\]

Nous allons surtout insister sur les pièges et les bonnes pratiques.

\subsection{Toujours utiliser les versions avec \texttt{n}}

Beaucoup de fonctions existent en deux versions :

\begin{itemize}
    \item \texttt{strcpy} / \texttt{strncpy}
    \item \texttt{strcat} / \texttt{strncat}
    \item \texttt{strcmp} / \texttt{strncmp}
\end{itemize}

\[
\boxed{\text{Recommandation : toujours utiliser les versions avec } n}
\]

Pourquoi ?

Parce qu’elles permettent de contrôler explicitement
la taille maximale manipulée.

\subsection{Copie : \texttt{strncpy}}

Prototype :

\begin{lstlisting}[language=C]
char* strncpy(char* dest, const char* src, size_t n);
\end{lstlisting}

\begin{itemize}
    \item \texttt{dest} : destination (modifiable)
    \item \texttt{src} : source (const)
    \item \texttt{n} : nombre maximal de caractères copiés
\end{itemize}

\textbf{Attention importante :}

\begin{itemize}
    \item Si la longueur de \texttt{src} ≥ n,
    \item alors \texttt{strncpy} ne garantit PAS l’ajout du caractère '\textbackslash 0'.
\end{itemize}

C’est un piège classique.

D’où la recommandation :

\[
\boxed{\text{Toujours allouer } taille + 1}
\]

et travailler avec une variable \texttt{taille}
qui représente uniquement les caractères utiles.

\subsection{Concaténation : \texttt{strncat}}

\begin{lstlisting}[language=C]
char* strncat(char* dest, const char* src, size_t n);
\end{lstlisting}

Elle concatène au plus n caractères de \texttt{src}
à la fin de \texttt{dest}.

Contrairement à \texttt{strncpy},
\texttt{strncat} garantit la terminaison par '\textbackslash 0'.

Toujours vérifier la taille disponible avant concaténation.

\subsection{Comparaison : \texttt{strncmp}}

\begin{lstlisting}[language=C]
int strncmp(const char* s1, const char* s2, size_t n);
\end{lstlisting}

Retour :

\begin{itemize}
    \item < 0 si s1 < s2 (ordre lexicographique)
    \item = 0 si égalité
    \item > 0 si s1 > s2
\end{itemize}

\textbf{Erreur classique :}

\begin{lstlisting}[language=C]
if (s1 == s2)   // FAUX
\end{lstlisting}

Cela compare les adresses,
pas le contenu.

La bonne façon :

\begin{lstlisting}[language=C]
if (strncmp(s1, s2, n) == 0) {
    // egalite
}
\end{lstlisting}

Souvent on écrit :

\begin{lstlisting}[language=C]
if (!strncmp(s1, s2, n)) {
    // egalite
}
\end{lstlisting}

\subsection{Longueur : \texttt{strlen}}

\begin{lstlisting}[language=C]
size_t strlen(const char* s);
\end{lstlisting}

Important :

\begin{itemize}
    \item Une chaîne ne connaît pas sa longueur.
    \item \texttt{strlen} parcourt les caractères jusqu’à '\textbackslash 0'.
\end{itemize}

Complexité :

\[
\boxed{O(n)}
\]

Si vous appelez \texttt{strlen} dans une boucle,
vous pouvez transformer un algorithme linéaire en quadratique.

Toujours stocker la longueur si elle est utilisée fréquemment.

\subsection{Recherche de caractères}

\paragraph{\texttt{strchr}}

\begin{lstlisting}[language=C]
char* strchr(const char* s, int c);
\end{lstlisting}

Retourne un pointeur vers la première occurrence de \texttt{c}
dans \texttt{s}, ou \texttt{NULL}.

\paragraph{\texttt{strrchr}}

Même chose, mais recherche de droite à gauche.

\subsection{Recherche de sous-chaîne : \texttt{strstr}}

\begin{lstlisting}[language=C]
char* strstr(const char* s1, const char* s2);
\end{lstlisting}

Retourne un pointeur vers la première occurrence de
la sous-chaîne \texttt{s2} dans \texttt{s1}.

\subsection{Entrées / sorties}

Affichage

\begin{lstlisting}[language=C]
printf("%s", s);
\end{lstlisting}

puts

\begin{lstlisting}[language=C]
puts(s);
\end{lstlisting}

Ajoute automatiquement un retour à la ligne.

fputs

\begin{lstlisting}[language=C]
fputs(s, stdout);
\end{lstlisting}

N’ajoute PAS de retour à la ligne.

\subsection{Lecture}

scanf

\begin{lstlisting}[language=C]
scanf("%s", s);
\end{lstlisting}

Attention :

\begin{itemize}
    \item Ne pas mettre \texttt{\&s}
    \item \texttt{s} est déjà un pointeur
\end{itemize}

Mais \texttt{scanf} ne protège pas contre les débordements.

fgets (recommandé)

\begin{lstlisting}[language=C]
fgets(s, taille, stdin);
\end{lstlisting}

\begin{itemize}
    \item Protège contre les dépassements
    \item Lit au maximum \texttt{taille - 1} caractères
    \item Ajoute toujours '\textbackslash 0'
\end{itemize}

\section{Résumé des bonnes pratiques}

\begin{itemize}
    \item Toujours allouer \texttt{n + 1}.
    \item Toujours utiliser les versions avec \texttt{n}.
    \item Ne jamais comparer deux chaînes avec \texttt{==}.
    \item Toujours vérifier les tailles.
    \item Éviter \texttt{scanf} non contrôlé.
    \item Stocker la longueur si utilisée souvent.
\end{itemize}

\[
\boxed{
\text{En C, une chaîne est fragile : elle repose uniquement sur '\textbackslash 0'.}
}
\]

\section{Pointeurs sur fonctions et généricité}

En C, on peut pointer vers n’importe quelle zone mémoire,
y compris vers du code.

Cela signifie :

\[
\boxed{\text{On peut avoir des pointeurs sur fonctions.}}
\]

\subsection{Déclaration d’un pointeur sur fonction}

Supposons la fonction :

\begin{lstlisting}[language=C]
double f(int x);
\end{lstlisting}

Pour déclarer un pointeur vers cette fonction :

\begin{lstlisting}[language=C]
double (*g)(int);
\end{lstlisting}

Règle :

\[
\text{On remplace le nom de la fonction par } (*nom\_pointeur)
\]

\subsection{Affectation}

\begin{lstlisting}[language=C]
g = &f;   // syntaxe pedagogique
\end{lstlisting}

On peut aussi écrire :

\begin{lstlisting}[language=C]
g = f;    // equivalent
\end{lstlisting}

En C, le nom d’une fonction est déjà un pointeur vers la fonction.

\subsection{Appel via pointeur}

Version explicite :

\begin{lstlisting}[language=C]
double y = (*g)(10);
\end{lstlisting}

Version simplifiée :

\begin{lstlisting}[language=C]
double y = g(10);
\end{lstlisting}

Les deux sont équivalentes.

\section{Passage de fonctions en argument}

Exemple classique : intégration numérique.

\begin{lstlisting}[language=C]
typedef double (*fonction)(double);

double integrer(fonction f, double a, double b);
\end{lstlisting}

Appel :

\begin{lstlisting}[language=C]
double resultat = integrer(sin, 0.0, M_PI);
\end{lstlisting}

On passe ici un pointeur sur la fonction \texttt{sin}.

\section{Généricité et pointeurs génériques}

Le type :

\begin{lstlisting}[language=C]
void*
\end{lstlisting}

ne signifie pas « pointeur vers rien ».

Il signifie :

\[
\boxed{\text{pointeur vers quelque chose dont on ignore le type}}
\]

C’est un pointeur générique.

\section{Exemple : tri générique avec qsort}

Prototype standard :

\begin{lstlisting}[language=C]
void qsort(void* base,
           size_t nmemb,
           size_t size,
           int (*compar)(const void*, const void*));
\end{lstlisting}

Arguments :

\begin{itemize}
\item \texttt{base} : pointeur vers le tableau
\item \texttt{nmemb} : nombre d’éléments
\item \texttt{size} : taille d’un élément
\item \texttt{compar} : fonction de comparaison
\end{itemize}

\subsection{Pourquoi void* ?}

Parce que le tableau peut contenir :

\begin{itemize}
\item des int,
\item des double,
\item des structures,
\item etc.
\end{itemize}

On ne connaît pas le type à l’avance.

\subsection{Exemple : tri d’un tableau d’entiers}

\begin{lstlisting}[language=C]
int tab[] = {4, 1, 9, 2};
size_t n = 4;

qsort(tab, n, sizeof(int), compare_int);
\end{lstlisting}

\subsection{Fonction de comparaison générique}

Prototype imposé :

\begin{lstlisting}[language=C]
int compare_int(const void* a, const void* b)
\end{lstlisting}

Implémentation propre :

\begin{lstlisting}[language=C]
int compare_int(const void* a, const void* b) {
    const int* i = a;
    const int* j = b;

    if (*i < *j) return -1;
    if (*i > *j) return 1;
    return 0;
}
\end{lstlisting}

\subsection{Pourquoi faire cette conversion ?}

Parce que :

\begin{itemize}
\item \texttt{a} et \texttt{b} sont des \texttt{void*}
\item On doit les interpréter comme des \texttt{int*}
\item On ne modifie pas les valeurs → \texttt{const}
\end{itemize}

\section{Comparaison pour des doubles}

\begin{lstlisting}[language=C]
int compare_double(const void* a, const void* b) {
    const double* x = a;
    const double* y = b;

    if (*x < *y) return -1;
    if (*x > *y) return 1;
    return 0;
}
\end{lstlisting}

Même schéma, type différent.

\section{Principe fondamental de la généricité en C}

En C, la généricité repose sur :

\begin{itemize}
\item \texttt{void*} pour les données,
\item pointeurs sur fonctions pour les comportements.
\end{itemize}

Contrairement aux templates C++,
tout est fait manuellement.

\[
\boxed{
\text{void* + pointeur sur fonction = généricité en C}
}
\]

\section{Points importants}

\begin{itemize}
\item Toujours utiliser \texttt{const void*} dans les comparateurs.
\item Toujours convertir vers le bon type avant déréférencement.
\item Ne jamais modifier les données dans une fonction de comparaison.
\item Toujours passer \texttt{sizeof(type)} à \texttt{qsort}.
\end{itemize}

\section{Lecture des types complexes en C}

À ce stade, nous avons vu :

\begin{itemize}
\item des pointeurs sur int,
\item des pointeurs sur pointeurs,
\item des fonctions retournant des pointeurs,
\item des pointeurs sur fonctions,
\item des tableaux de pointeurs,
\item des combinaisons de tout cela.
\end{itemize}

Il devient alors essentiel de savoir lire des déclarations complexes.

\subsection{Règle fondamentale de lecture}

\[
\boxed{\text{Toujours commencer par le nom.}}
\]

Ensuite :

\begin{itemize}
\item Respecter les parenthèses.
\item Lire d’abord ce qui est immédiatement autour du nom.
\item Puis élargir progressivement vers l’extérieur.
\end{itemize}

Les parenthèses sont cruciales :
elles définissent les priorités.

\section{Exemple de type complexe}

Considérons la déclaration suivante :

\begin{lstlisting}[language=C]
int *(*(*f())[])();
\end{lstlisting}

Comment lire cela ?

\subsection{Étape 1 : partir du nom}

Le nom est :

\[
f
\]

On regarde immédiatement ce qui suit :

\begin{lstlisting}
f()
\end{lstlisting}

Donc :

\[
f \text{ est une fonction}
\]

\subsection{Étape 2 : que retourne f ?}

On regarde à gauche de \texttt{f()} :

\begin{lstlisting}
(*f())
\end{lstlisting}

Donc :

\[
f \text{ retourne un pointeur}
\]

\subsection{Étape 3 : vers quoi pointe ce pointeur ?}

On regarde :

\begin{lstlisting}
(*(*f())[])
\end{lstlisting}

Donc :

\[
f \text{ retourne un pointeur vers un tableau}
\]

\subsection{Étape 4 : contenu du tableau}

On regarde :

\begin{lstlisting}
(*(*f())[])()
\end{lstlisting}

Cela signifie :

\[
\text{tableau de pointeurs sur fonctions}
\]

\subsection{Étape 5 : que retournent ces fonctions ?}

On regarde le type final :

\begin{lstlisting}
int *
\end{lstlisting}

Donc :

\[
\text{les fonctions retournent un pointeur sur int}
\]

\section{Lecture complète}

On peut lire la déclaration ainsi : \\

f est une fonction qui retourne un pointeur vers un tableau de pointeurs sur fonctions retournant des pointeurs sur int

\section{Conseil pratique essentiel}

Même si vous êtes capable de lire ce genre de déclaration :

\[
\boxed{\text{N’écrivez jamais cela directement.}}
\]

Utilisez des \texttt{typedef}.

\section{Réécriture avec typedef}

\begin{lstlisting}[language=C]
typedef int* int_ptr;

typedef int_ptr (*func_ptr)(void);

typedef func_ptr func_array[];

typedef func_array* return_type;

return_type f(void);
\end{lstlisting}

Cette version est :

\begin{itemize}
\item plus lisible,
\item plus maintenable,
\item moins sujette aux erreurs.
\end{itemize}

\section{Casting en C : valeurs et pointeurs}

En C, le typage est souple.  
On peut convertir des expressions vers un autre type via un \textbf{casting}.

Syntaxe générale :

\begin{lstlisting}[language=C]
(type) expression
\end{lstlisting}

Il est essentiel de distinguer :

\begin{itemize}
\item le casting de valeurs,
\item le casting de pointeurs.
\end{itemize}

\section{Casting de valeurs}

Exemple :

\begin{lstlisting}[language=C]
double x = 5.4;
int y = (int)x;
\end{lstlisting}

Ici :

\begin{itemize}
\item x est une valeur double.
\item (int)x convertit la valeur.
\item La partie fractionnaire est supprimée.
\end{itemize}

Donc :

\[
y = 5
\]

Le comportement est un arrondi vers zéro :

\begin{itemize}
\item 5.9 → 5
\item -5.9 → -5
\end{itemize}

Ici, il y a réellement transformation de valeur.

\section{Casting de pointeurs}

C’est ici que les choses deviennent délicates.

Considérons :

\begin{lstlisting}[language=C]
double x = 5.4;
int* i = (int*)&x;
\end{lstlisting}

Important :

\[
\boxed{
\text{Le casting de pointeur ne modifie pas la mémoire.}
}
\]

Il change uniquement \textbf{l’interprétation} des bits.

\subsection{Ce qui se passe réellement}

En mémoire :

\begin{itemize}
\item x occupe sizeof(double) octets.
\item Ces octets représentent 5.4 en format IEEE.
\item i pointe vers la même adresse.
\item Mais i lit seulement sizeof(int) octets.
\end{itemize}

Donc :

\begin{lstlisting}[language=C]
printf("%d\n", *i);
\end{lstlisting}

affichera une valeur incohérente correspondant
à l’interprétation binaire des premiers octets du double.

Il n’y a aucune conversion de valeur.

\section{Différence fondamentale}

\begin{itemize}
\item Casting de valeur → conversion réelle.
\item Casting de pointeur → relecture brute des bits.
\end{itemize}

\[
\boxed{
\text{Un casting de pointeur change la lecture, pas les données.}
}
\]

\section{Casting et code générique}

Le casting est utilisé légitimement dans :

\begin{itemize}
\item les fonctions génériques,
\item les manipulations de void*,
\item qsort.
\end{itemize}

Exemple :

\begin{lstlisting}[language=C]
int compare_int(const void* a, const void* b) {
    const int* i = a;
    const int* j = b;

    if (*i < *j) return -1;
    if (*i > *j) return 1;
    return 0;
}
\end{lstlisting}

Ici, le cast est nécessaire
pour interpréter correctement la mémoire.

\section{Casting de pointeurs sur fonctions}

On peut aussi caster des pointeurs sur fonctions.

Exemple avec \texttt{qsort} :

\begin{lstlisting}[language=C]
qsort(tab, n, sizeof(Person),
      (int (*)(const void*, const void*))compare_person);
\end{lstlisting}

Cependant, cette écriture est difficile à lire.

Une alternative plus propre :

\begin{lstlisting}[language=C]
int compare_person(const void* a, const void* b) {
    const Person* p1 = a;
    const Person* p2 = b;

    ...
}
\end{lstlisting}

Puis :

\begin{lstlisting}[language=C]
qsort(tab, n, sizeof(Person), compare_person);
\end{lstlisting}

Plus lisible et plus sûr.

\section{Casting et void*}

En C :

\begin{itemize}
\item On peut convertir automatiquement vers void*.
\item On peut convertir automatiquement depuis void*.
\end{itemize}

Exemple :

\begin{lstlisting}[language=C]
int* p1;
void* p2;

p2 = p1;   // pas de cast necessaire
p1 = p2;   // pas de cast necessaire en C
\end{lstlisting}

\subsection{Erreur courante avec malloc}

Beaucoup écrivent :

\begin{lstlisting}[language=C]
int* p = (int*)malloc(sizeof(int));
\end{lstlisting}

Ce cast est inutile en C.

Écriture correcte :

\begin{lstlisting}[language=C]
int* p = malloc(sizeof(int));
\end{lstlisting}

Pourquoi ?

\begin{itemize}
\item malloc retourne un void*.
\item La conversion vers int* est implicite en C.
\end{itemize}

Ajouter un cast :

\begin{itemize}
\item surcharge le code,
\item peut masquer des erreurs d’inclusion de headers.
\end{itemize}
